\scanpage{154}
\mspage{70}
% archive sheet number: 153

\uline{Dém} \deleted{\unsure{Il}} suffit de \deleted{\unsure{noter|voir}} \inserted{noter} que si $S,\alpha$ est comme dans 4.3.4.\deleted{\unsure{a}}b), alors $\alpha$ est accessible \unsure{sss|ssi} \unsure{les|le} $\rho(S)$ \deleted{\unsure{$(s\in\mathrm{Ob}$}} est formé d'objets accessibles de $M$ ; comme cette partie est strict\textsuperscript{mt} génératrice, on sait que cela implique que \uline{tt} objet de $M$ est accessible. (Cf SGA 4 I 9.9.)

\deleted{\uline{\unsure{Corollaire}} \unsure{4.3.6}}

Je vais essayer de donner un énoncé plus complet :

\uline{Théorème} 4.3.6. Soit $M$ une $\unsure{U|\mathcal U}$-catégorie Conditions équivalentes

a) $M$ stable par petites $\underrightarrow{\lim}$, et elle admet une petite sous-catégorie pleine, \inserted{$S$} strictement génératrice \unsure{$S$} et formée d'objets accessibles de $M$

a$'$) $M$ est une \inserted{catégorie $U$-}pseudo topos \inserted{\deleted{\unsure{faiblement}}} accessible, i.e. satisfait a), avec en plus \uline{\deleted{\unsure{tout|tt}}} élément de $M$ (pas seulement ceux de $S$) accessible dans $M$

b) Il existe une petite catégorie $S$, et un foncteur $\alpha:M\unsure{\lhook\joinrel-|\hookrightarrow} S^\wedge$ \uline{pleinement fidèle} \uline{et accessible}, tel que $\alpha$ admette un adjoint à gauche.

